\scanpage{23}
au \unsure{fait|cas|fond} où $\mathcal C$ serait la formalisation de la catégorie triviale définie par $X_0/\mathrm{Spec}\,Z$). En particulier
\[
\mathscr T^{\mathcal C}=\deleted{\unsure{\mathrm{dual}|\text{déf}}}\ \underline{\mathrm{Hom}}_{\underline O_{X_0}}(\Omega_{\mathcal C}^{\deleted{\unsure{\mathcal C}}},\underline O_X)=\varinjlim_{\alpha}\underline{\mathrm{Hom}}_{\underline O_{X_0}}(\Omega_{\unsure{G_\alpha|\mathcal C_\alpha}}^{\mathcal C},\underline O_{X_0})
\]
opère sur $\underline O_{X_0}$ par \uline{dérivations}\,\ldots. Cf aussi plus bas pour sa déduction à partir du complexe de De Rham $\Omega_{\mathcal C}^{*\deleted{\unsure{\mathcal C}}}$ de $\mathcal C$.

2.\quad \uline{Complexes de De Rham} \uline{généraux}.

Appelons complexe de De Rham construit sur un \deleted{\unsure{sché}} Module \inserted{qu. ch.} $\Omega$ sur un schéma $X_0$ [NB marchera aussi pour tout espace annelé, de même que la suite du 1.] \inserted{tout complexe obtenu par} une \deleted{\unsure{op}} différentielle d'anneau gradué sur $\bigwedge^{*}\Omega$.

\uline{Remarque} 1\quad Cette différentielle est connue quand on la connaît en degrés 0 et 1
\[
\begin{aligned}
\delta_0:\underline O_{X_0}\longrightarrow\Omega\quad &\inserted{\alpha)\ \text{\uline{dérivation}}}\quad [\,\text{i.e.}\quad\Omega^1_{X_0/\mathbb Z}\xrightarrow{u}\Omega\\
\inserted{\alpha)\ \delta_0(fg)=\delta_0(f)g+f\delta_0(g)}\quad &\text{par}\quad\delta(f)=u(df)\quad ]\\
\delta_1:\Omega\longrightarrow\bigwedge^2\Omega &
\end{aligned}
\]
Condition nécessaire sur $\delta_1$
\[
\left\{\begin{aligned}
\beta)\ &\delta_1(f\omega)=\delta_0(f)\,\delta_1(\omega)+f\,\delta_1(\omega)\\
\gamma)\ &\delta_1\delta_0=0
\end{aligned}\right.
\]

De plus, on peut choisir $\delta_0$ et $\delta_1$ arbitrairement, de façon à satisfaire ces conditions : ils se prolongent (de façon

% archive sheet number: 22

